% ══════════════════════════════════════════════════════════════════════════════
%  BIODIVERSITY AND ECOSYSTEMS — A5 LESSON BOOKLET, LESSONS 1–4
%  Written for the 2027 workbook style from Bio_Ecosystems_Unit_Booklet_Yr7_8.tex
%  (same ideas, tasks and First Nations content; restructured into Parts A–D).
%  Built by make_a5.py. Lesson numbers = the unit's planner (planner-data.json).
% ══════════════════════════════════════════════════════════════════════════════

% ══════════════════════════════════════════════════════════════════════════════
%  LESSON 1 — LIVING OR NON-LIVING?
% ══════════════════════════════════════════════════════════════════════════════
\setcounter{section}{0}
\sectiondiv{Living or Non-living?}{MRS GREN \textbullet{} The seven life processes \textbullet{} Dead versus never alive}\label{sec:1}
\lessoncodes{VC2S8U01 \textbullet{} VC2S8H02}

\vocabhead
\begin{ctab}
\begin{tabular}{|SW{3.4cm}|W{7.0cm}|}
  \hline\hd & \thead{Word} & \thead{Definition (copy from the slide)} \\ \hline
  & \scaf{\textbf{Organism}} & \scaf{A single living thing, such as a plant, an animal, a fungus or a bacterium.} \\ \hline
  & \textbf{Life process} & \\ \hline
  & \textbf{Respiration} & \\ \hline
  & \textbf{Excretion} & \\ \hline
  & \textbf{Sensitivity} & \\ \hline
  & \textbf{Cell} & \\ \hline
\end{tabular}
\end{ctab}

% ─── PART A ───────────────────────────────────────────────────────────────────
\lessonpart{A}{Reading and comprehension}{Teal}\label{read:1}

\begin{readingbox}{Reading --- What Makes Something Alive?}
A kangaroo is alive. A rock is not. That much is easy. But what about a seed that has sat in a packet for five years, a mushroom, or a flame that seems to eat, breathe and grow?

Scientists decide whether an \textbf{organism} is alive by checking for seven \textbf{life processes}. Every living thing carries out all seven, at least at some point in its life. They are remembered with the name \textbf{MRS GREN}: Movement, Respiration, Sensitivity, Growth, Reproduction, Excretion and Nutrition.

Some of these are easy to see. A possum moves, grows and has young. Others happen out of sight. \textbf{Respiration} is a chemical reaction in every cell that releases energy from food. \textbf{Excretion} is getting rid of the waste that the body's chemical reactions make. \textbf{Sensitivity} is detecting changes around you and responding to them: a plant bends towards light, and a lizard runs from a shadow.

There is one more test. Every living thing is made of one or more \textbf{cells}, the tiny building blocks of life. A thing that was once alive, like a dry leaf or a wooden desk, is \emph{dead}. A thing that was never alive, like a river rock, is \emph{non-living}.
\end{readingbox}

\rcreset
\rcq{What does each letter of MRS GREN stand for?}{2.5cm}
\rcq{What is respiration, and where does it happen?}{2.5cm}
\rcq{Give an example of sensitivity in a plant and in an animal.}{2.5cm}
\rcq{What is the difference between \emph{dead} and \emph{non-living}?}{2.5cm}
\rcq{What are all living things made of?}{1.5cm}

% ─── PART B ───────────────────────────────────────────────────────────────────
\lessonpart{B}{Theory and examples}{Blue}\label{theory:1}

\subsection{The Seven Life Processes}\label{sub:1.1}

\begin{center}
\begin{tikzpicture}[font=\small,
  lp/.style={draw=Ink, line width=0.7pt, rounded corners=1.2mm, fill=white, minimum width=2.5cm, minimum height=0.75cm, align=center}]
  \node[draw=Ink, line width=1.2pt, circle, minimum size=1.9cm, font=\headingfont\bfseries, align=center] (c) at (0,0) {MRS\\GREN};
  \foreach \a/\t [count=\i] in {90/{\textbf{M}ovement}, 38.6/{\textbf{R}espiration}, -12.9/{\textbf{S}ensitivity}, -64.3/{\textbf{G}rowth}, -115.7/{\textbf{R}eproduction}, -167.1/{\textbf{E}xcretion}, 141.4/{\textbf{N}utrition}}{
    \node[lp] (n\i) at (\a:2.75) {\t};
    \draw[Ink, line width=0.6pt] (c) -- (n\i);}
\end{tikzpicture}
\figcap{The seven life processes.}
\end{center}

\begin{skillbox}{MRS GREN}
\renewcommand{\arraystretch}{1.15}
\begin{tabular}{@{}p{2.4cm}p{7.3cm}@{}}
\textbf{Movement} & All living things move in some way; even plants move slowly towards light. \\
\textbf{Respiration} & Releasing energy from food, inside every cell. \\
\textbf{Sensitivity} & Detecting changes in the surroundings and responding to them. \\
\textbf{Growth} & A lasting increase in size or in the number of cells. \\
\textbf{Reproduction} & Making new organisms (offspring). \\
\textbf{Excretion} & Getting rid of waste made by the body's chemical reactions. \\
\textbf{Nutrition} & Taking in and using food for energy and for building the body. \\
\end{tabular}
\end{skillbox}

\begin{workedex}{Is fire alive?}
Fire spreads (movement), uses oxygen (like respiration), responds to wind (sensitivity) and gets bigger (growth).

\smallskip
But fire does not make offspring of its own kind from cells, it does not excrete waste from a body, and it does not use food to build a body. Most importantly, it is \textbf{not made of cells}.

\smallskip
\textbf{Conclusion:} fire is \textbf{non-living}. It shows some of MRS GREN, but not all seven, and it has no cells.
\end{workedex}

\begin{wedo}{Is it alive?}
A crystal of salt grows larger when it sits in a salty solution. Use MRS GREN and the cell test to decide whether the crystal is alive.
\answerbox{5cm}
\end{wedo}

\subsection{Living, Dead or Non-living?}\label{sub:1.2}

\begin{skillbox}{Three groups}
\textbf{Living:} carries out all seven life processes now. \quad \textbf{Dead:} was once living, made of cells, but no longer carries out the life processes (a dry leaf, a wooden desk, a bone). \quad \textbf{Non-living:} was never alive and is not made of cells (a rock, water, a candle flame).
\end{skillbox}

\begin{workedex}{Classifying a seed}
A wattle seed has sat in a paper bag for three years. It does not move, grow or feed. Is it dead?

\smallskip
\textbf{No.} A seed is \textbf{living}, but dormant: its cells respire very slowly. Put it in warm, moist soil and it germinates, grows and later reproduces. It only looks dead because the life processes are happening slowly.
\end{workedex}

\begin{wedo}{Classifying a seed}
A cooked pea and a dried pea from a seed packet look alike. Which one is living, and how could you test it?
\answerbox{5cm}
\end{wedo}

\begin{countrybox}{First Nations perspectives --- Country as a living system}
Wurundjeri Woi-wurrung people understand Country --- land, water, air, plants, animals and spirit --- as one living, connected system, not a set of separate resources. When Elders say Country is ``sick'' or ``healthy'', they describe real ecological states, built on tens of thousands of years of careful observation. MRS GREN describes one organism at a time; Country describes how all of them live together.
\end{countrybox}

% ─── PART C ───────────────────────────────────────────────────────────────────
\lessonpart{C}{Practice}{Amber}\label{prac:1}

\task Sort each item into the right group with a tick, then give a reason.
\begin{ctab}
\begin{tabular}{|SW{2.4cm}|Z{1.2cm}|Z{1.2cm}|Z{1.3cm}|W{3.3cm}|}
  \hline\hd & \thead{Item} & \thead{Living} & \thead{Dead} & \thead{Non-living} & \thead{Reason} \\ \hline
  & \scaf{Dry leaf} & & \scaf{\ding{51}} & & \scaf{Made of cells, was alive, no life processes now} \\ \hline
  & Candle flame & & & & \\ \hline
  & Bacterium & & & & \\ \hline
  & Wooden desk & & & & \\ \hline
  & Yeast in bread dough & & & & \\ \hline
  & River rock & & & & \\ \hline
  & Mushroom & & & & \\ \hline
\end{tabular}
\end{ctab}

\needspace{10\baselineskip}
\task A virus can make copies of itself, but only inside another organism's cell. On its own it does not grow, move, feed or respire, and it is not made of cells. Is a virus living or non-living? Explain using MRS GREN.
\answerbox{5cm}

\needspace{10\baselineskip}
\task A car moves, uses fuel, gives out waste gases and has sensors that respond to the road. Explain why a car is still non-living.
\answerbox{5cm}

\needspace{10\baselineskip}
\task Breathing is not the same as respiration. Explain the difference.
\answerbox{5cm}

% ─── PART D ───────────────────────────────────────────────────────────────────
\lessonpart{D}{Above level}{Violet}\label{ext:1}

\begin{yr8box}
These questions go beyond the Year 7 standard.

\medskip
\begin{keepext}
\exttask Plant cells and animal cells both carry out all seven life processes. Describe two parts of a plant cell that an animal cell does not have, and explain how one of them lets a plant make its own food.
\answerbox{5cm}
\end{keepext}
\begin{keepext}
\exttask Some scientists argue that viruses should count as living, because they change over time and carry genetic material. Give one piece of evidence for and one against this argument, and state your own view.
\answerbox{5cm}
\end{keepext}
\end{yr8box}

% ══════════════════════════════════════════════════════════════════════════════
%  LESSON 2 — CLASSIFICATION: LINNAEAN LEVELS
% ══════════════════════════════════════════════════════════════════════════════
\setcounter{section}{1}
\sectiondiv{Classification: Linnaean Levels}{Why we classify \textbullet{} Kingdom to species \textbullet{} What classification tells us}\label{sec:2}
\lessoncodes{VC2S8U01 \textbullet{} VC2S8H02}

\vocabhead
\begin{ctab}
\begin{tabular}{|SW{3.4cm}|W{7.0cm}|}
  \hline\hd & \thead{Word} & \thead{Definition (copy from the slide)} \\ \hline
  & \scaf{\textbf{Classification}} & \scaf{Sorting living things into groups by the features they share.} \\ \hline
  & \textbf{Taxonomy} & \\ \hline
  & \textbf{Taxon} & \\ \hline
  & \textbf{Kingdom} & \\ \hline
  & \textbf{Species} & \\ \hline
  & \textbf{Vertebrate} & \\ \hline
\end{tabular}
\end{ctab}

% ─── PART A ───────────────────────────────────────────────────────────────────
\lessonpart{A}{Reading and comprehension}{Teal}\label{read:2}

\begin{readingbox}{Reading --- A Place for Every Living Thing}
Scientists have named about two million species, and millions more are still waiting to be found. To make sense of so many, they use \textbf{classification}: sorting living things into groups by the features they share.

In the 1700s the Swedish scientist Carl Linnaeus built the system we still use. The science of naming and classifying organisms is called \textbf{taxonomy}. Linnaeus placed every organism into a set of groups, one inside the next, like addresses: country, state, suburb, street, house. Each group is a \textbf{taxon}.

The biggest group is the \textbf{kingdom}. All animals are in one kingdom, all plants in another. Each kingdom is split into smaller and smaller groups: phylum, class, order, family and genus. The smallest group is the \textbf{species}: organisms so alike that they can breed together and have young that can also breed.

The deeper two organisms sit together, the more closely related they are. A kangaroo and a wombat share every level down to the order. A kangaroo and a snake share only the kingdom and the phylum: both are \textbf{vertebrates}, animals with a backbone.
\end{readingbox}

\rcreset
\rcq{Why do scientists classify living things?}{2.5cm}
\rcq{Who built the classification system we still use, and when?}{2cm}
\rcq{What is the biggest group, and what is the smallest?}{2cm}
\rcq{What makes two organisms members of the same species?}{2.5cm}
\rcq{How can you tell from the groups that two organisms are closely related?}{2.5cm}

% ─── PART B ───────────────────────────────────────────────────────────────────
\lessonpart{B}{Theory and examples}{Blue}\label{theory:2}

\subsection{The Seven Levels}\label{sub:2.1}

\begin{center}
\begin{tikzpicture}[font=\small]
  \foreach \t [count=\i from 0] in {Kingdom,Phylum,Class,Order,Family,Genus,Species}{
    \pgfmathsetmacro\w{4.0-0.5*\i}\pgfmathsetmacro\y{-0.62*\i}
    \pgfmathsetmacro\g{4+4*\i}
    \draw[Ink, line width=0.7pt, fill=black!\g] (-\w,\y) -- (\w,\y) -- (\w-0.5,\y-0.55) -- (-\w+0.5,\y-0.55) -- cycle;
    \node[font=\small\bfseries] at (0,\y-0.275) {\t};}
  \draw[Ink, line width=0.9pt, -{Stealth[length=5pt]}] (4.5,0) -- (4.5,-4.3);
  \node[font=\scriptsize, anchor=west, align=left] at (4.6,-0.2) {broadest};
  \node[font=\scriptsize, anchor=west, align=left] at (4.6,-4.1) {most\\specific};
\end{tikzpicture}
\figcap{The seven Linnaean levels.}
\end{center}

\begin{skillbox}{Remembering the order}
\textbf{K}ing \textbf{P}hillip \textbf{C}ame \textbf{O}ver \textbf{F}or \textbf{G}ood \textbf{S}paghetti: Kingdom, Phylum, Class, Order, Family, Genus, Species. Each level down holds fewer kinds of organism, and they share more features.
\end{skillbox}

\begin{workedex}{Classifying the eastern grey kangaroo}
\renewcommand{\arraystretch}{1.15}
\begin{tabular}{@{}p{1.6cm}p{2.5cm}p{5.6cm}@{}}
\textbf{Level} & \textbf{Group} & \textbf{What it tells us} \\ \hline
Kingdom & Animalia & An animal: many cells, eats its food \\
Phylum & Chordata & Has a backbone (a vertebrate) \\
Class & Mammalia & Has fur and feeds its young on milk \\
Order & Diprotodontia & Two front teeth in the lower jaw; with wombats and possums \\
Family & Macropodidae & ``Big feet'': kangaroos and wallabies \\
Genus & \textit{Macropus} & Large kangaroos \\
Species & \textit{giganteus} & This one species: ``giant'' \\
\end{tabular}
\end{workedex}

\begin{wedo}{Comparing two classifications}
The common wombat is Animalia, Chordata, Mammalia, Diprotodontia, Vombatidae, \textit{Vombatus}, \textit{ursinus}. Compare it with the eastern grey kangaroo above. At which level do they first differ, and what does that tell us about how closely they are related?
\answerbox{5cm}
\end{wedo}

\begin{countrybox}{First Nations perspectives --- Other ways to classify}
Wurundjeri Woi-wurrung people have their own detailed systems for classifying plants and animals, built over tens of thousands of years. They group living things not only by how they look but by what they do: what they eat, where they shelter, when they are active, and what they signal about the season. Certain bird calls, for example, mean a plant is ready to harvest or fish are running. This is not ``less scientific'' than the Linnaean system --- it is organised for a different purpose.
\end{countrybox}

% ─── PART C ───────────────────────────────────────────────────────────────────
\lessonpart{C}{Practice}{Amber}\label{prac:2}

\task Sort these six animals into groups using any feature you can observe: eastern grey kangaroo, common wombat, tiger snake, wedge-tailed eagle, platypus, southern brown bandicoot. Draw your groups and label each group with the feature you used.
\drawbox{9cm}{}

\needspace{26\baselineskip}
\task The short-beaked echidna's groups, from kingdom to species, are: Animalia, Chordata, Mammalia, Monotremata, Tachyglossidae, \textit{Tachyglossus}, \textit{aculeatus}. Complete the table.
\begin{ctab}
\begin{tabular}{|SW{1.9cm}|W{2.8cm}|W{4.6cm}|}
  \hline\hd & \thead{Level} & \thead{Echidna} & \thead{What it tells us} \\ \hline
  & \scaf{Kingdom} & \scaf{Animalia} & \scaf{It is an animal} \\ \hline
  & Phylum & & \\ \hline
  & Class & & \\ \hline
  & Order & & \\ \hline
  & Family & & \\ \hline
  & Genus & & \\ \hline
  & Species & & \\ \hline
\end{tabular}
\end{ctab}

\needspace{10\baselineskip}
\task The eastern grey kangaroo's order is Diprotodontia and the echidna's is Monotremata; they share kingdom, phylum and class. At which level do they first differ? What does this tell us about how closely they are related?
\answerbox{5cm}

\needspace{10\baselineskip}
\task Why is it useful to sort organisms into a ladder of groups, instead of just giving each one its own name?
\answerbox{5cm}

% ─── PART D ───────────────────────────────────────────────────────────────────
\lessonpart{D}{Above level}{Violet}\label{ext:2}

\begin{yr8box}
These questions go beyond the Year 7 standard.

\medskip
\begin{keepext}
\exttask Scientists now place three \textbf{domains} above the kingdoms: Bacteria, Archaea and Eukarya. Explain what feature of the cells separates the domains, then place each of these in a domain: a human, the bacterium \textit{E.\ coli}, a microbe from a hot spring (\textit{Sulfolobus}).
\answerbox{5cm}
\end{keepext}
\begin{keepext}
\exttask Reading DNA has moved some organisms to different groups. Explain why DNA evidence is more reliable than appearance for working out how closely organisms are related.
\answerbox{5cm}
\end{keepext}
\end{yr8box}

% ══════════════════════════════════════════════════════════════════════════════
%  LESSON 3 — SCIENTIFIC NAMES AND THE SIX KINGDOMS
% ══════════════════════════════════════════════════════════════════════════════
\setcounter{section}{2}
\sectiondiv{Scientific Names and the Six Kingdoms}{Binomial names \textbullet{} The six kingdoms \textbullet{} Victorian examples}\label{sec:3}
\lessoncodes{VC2S8U01 \textbullet{} VC2S8H02}

\vocabhead
\begin{ctab}
\begin{tabular}{|SW{3.4cm}|W{7.0cm}|}
  \hline\hd & \thead{Word} & \thead{Definition (copy from the slide)} \\ \hline
  & \scaf{\textbf{Binomial name}} & \scaf{A species' two-part scientific name: genus, then species.} \\ \hline
  & \textbf{Genus} & \\ \hline
  & \textbf{Nucleus} & \\ \hline
  & \textbf{Eukaryote} & \\ \hline
  & \textbf{Prokaryote} & \\ \hline
  & \textbf{Photosynthesis} & \\ \hline
\end{tabular}
\end{ctab}

% ─── PART A ───────────────────────────────────────────────────────────────────
\lessonpart{A}{Reading and comprehension}{Teal}\label{read:3}

\begin{readingbox}{Reading --- One Animal, Many Names}
In Victoria a ``magpie'' is a black-and-white bird that swoops in spring. In England a ``magpie'' is a completely different bird. The same animal can have many common names, and the same common name can mean different animals. Scientists avoid the muddle with one name per species that is used in every country and language.

Each species has a two-part \textbf{binomial name}. The first part is the \textbf{genus}, with a capital letter. The second part is the species, all in lower case. The whole name is written in italics, or underlined when handwritten. The common wombat is \textit{Vombatus ursinus}: ``bear-like wombat''.

At the top of the ladder, every organism belongs to one of six kingdoms: animals, plants, fungi, protists, bacteria and archaea. The kingdoms are told apart by their cells. In animals, plants, fungi and protists, each cell has a \textbf{nucleus}, a control centre wrapped in a membrane: these are \textbf{eukaryotes}. Bacteria and archaea have no nucleus: they are \textbf{prokaryotes}.

How an organism gets its food matters too. Plants make their own food from sunlight by \textbf{photosynthesis}. Animals eat other organisms. Fungi, like mushrooms, soak up food from dead or living things around them.
\end{readingbox}

\rcreset
\rcq{Why do scientists need a scientific name as well as a common name?}{2.5cm}
\rcq{Give the two rules for writing the two parts of a binomial name.}{2.5cm}
\rcq{Name the six kingdoms.}{2cm}
\rcq{What is the difference between a eukaryote and a prokaryote?}{2.5cm}
\rcq{How do plants, animals and fungi each get their food?}{2.5cm}

% ─── PART B ───────────────────────────────────────────────────────────────────
\lessonpart{B}{Theory and examples}{Blue}\label{theory:3}

\subsection{Writing Scientific Names}\label{sub:3.1}

\begin{skillbox}{The rules}
\textbf{1.} Two words: genus, then species. \quad \textbf{2.} The genus starts with a capital letter. \quad \textbf{3.} The species is all lower case. \quad \textbf{4.} Italics when typed; underline each word when handwritten. \quad \textbf{5.} The same name is used all over the world.
\end{skillbox}

\begin{workedex}{Fixing a name}
\textbf{Wrong:} Ornithorhynchus Anatinus (the platypus).

\smallskip
The genus \emph{Ornithorhynchus} is right to have a capital. The species \emph{Anatinus} must be lower case. It is typed, so it must be in italics.

\smallskip
\textbf{Right:} \textit{Ornithorhynchus anatinus}.
\end{workedex}

\begin{wedo}{Fixing a name}
Rewrite this name correctly and say which rules were broken: \textsc{macropus giganteus} (the eastern grey kangaroo).
\answerbox{4cm}
\end{wedo}

\subsection{The Six Kingdoms}\label{sub:3.2}

\begin{center}
\begin{tikzpicture}[font=\small,
  k/.style={draw=Ink, line width=0.7pt, rounded corners=1.2mm, fill=white, text width=2.9cm, align=center, minimum height=1.45cm, inner sep=2pt}]
  \node[k] at (0,0) {\textbf{Animalia}\\{\scriptsize kangaroos, insects, fish}};
  \node[k] at (3.4,0) {\textbf{Plantae}\\{\scriptsize wattles, grasses, ferns}};
  \node[k] at (6.8,0) {\textbf{Fungi}\\{\scriptsize mushrooms, moulds, yeast}};
  \node[k] at (0,-1.85) {\textbf{Protista}\\{\scriptsize amoeba, algae}};
  \node[k, fill=black!8] at (3.4,-1.85) {\textbf{Bacteria}\\{\scriptsize \textit{E.\ coli}, \textit{Salmonella}}};
  \node[k, fill=black!8] at (6.8,-1.85) {\textbf{Archaea}\\{\scriptsize hot-spring microbes}};
  \node[font=\scriptsize, anchor=south] at (3.4,0.85) {Eukaryotes (cells with a nucleus): white};
  \node[font=\scriptsize, anchor=north] at (3.4,-2.7) {Prokaryotes (no nucleus): grey};
\end{tikzpicture}
\figcap{The six kingdoms.}
\end{center}

\begin{workedex}{Classifying into kingdoms}
\textbf{A mushroom in Sherbrooke Forest:} Fungi. It has many cells with a nucleus, and it cannot photosynthesise; it soaks up food from dead leaves and wood.\\
\textbf{A wattle tree:} Plantae. Its cells have a nucleus, and it photosynthesises.\\
\textbf{\textit{Salmonella}:} Bacteria. It is a single cell with no nucleus.
\end{workedex}

\begin{wedo}{Classifying into kingdoms}
Green pond scum is made of single cells, each with a nucleus, that photosynthesise. Which kingdom is it in? Give two reasons.
\answerbox{4cm}
\end{wedo}

\begin{countrybox}{First Nations perspectives --- Names that carry knowledge}
In Woi-wurrung, the language of the Wurundjeri people, the names of plants and animals often carry ecological knowledge: the kind of Country a plant grows on, the season an animal is active, or the insect that pollinates a flower. Scientific names carry information too, but they record relationships between species, while Woi-wurrung names record relationships with Country and people.
\end{countrybox}

% ─── PART C ───────────────────────────────────────────────────────────────────
\lessonpart{C}{Practice}{Amber}\label{prac:3}

\task Rewrite each scientific name correctly.
\begin{ctab}
\begin{tabular}{|SW{4.6cm}|W{5.4cm}|}
  \hline\hd & \thead{Written wrongly} & \thead{Written correctly} \\ \hline
  & \scaf{vombatus Ursinus} & \scaf{\textit{Vombatus ursinus}} \\ \hline
  & \textit{Macropus Giganteus} & \\ \hline
  & ornithorhynchus anatinus & \\ \hline
  & \textit{TACHYGLOSSUS aculeatus} & \\ \hline
\end{tabular}
\end{ctab}

\needspace{14\baselineskip}
\task Name the kingdom of each organism and give two reasons.
\begin{ctab}
\begin{tabular}{|SW{3.0cm}|W{2.2cm}|W{4.8cm}|}
  \hline\hd & \thead{Organism} & \thead{Kingdom} & \thead{Two reasons} \\ \hline
  & \scaf{Wattle tree} & \scaf{Plantae} & \scaf{Cells have a nucleus; it photosynthesises} \\ \hline
  & Amoeba (single cell with a nucleus) & & \\ \hline
  & Common wombat & & \\ \hline
  & Hot-spring microbe (no nucleus) & & \\ \hline
  & Bread mould & & \\ \hline
\end{tabular}
\end{ctab}

\needspace{10\baselineskip}
\task Why is it important that scientists all over the world use the same naming system? Give an example of the confusion common names can cause.
\answerbox{5cm}

% ─── PART D ───────────────────────────────────────────────────────────────────
\lessonpart{D}{Above level}{Violet}\label{ext:3}

\begin{yr8box}
These questions go beyond the Year 7 standard.

\medskip
\begin{keepext}
\exttask Fungi were once classified as plants. Describe two features of fungi that make them different from plants, and explain why these justify a separate kingdom.
\answerbox{5cm}
\end{keepext}
\begin{keepext}
\exttask Draw and label a simple prokaryotic cell and a simple eukaryotic cell, showing the key difference between them.
\drawbox{5.5cm}{}
\end{keepext}
\end{yr8box}

% ══════════════════════════════════════════════════════════════════════════════
%  LESSON 4 — DICHOTOMOUS KEYS
% ══════════════════════════════════════════════════════════════════════════════
\setcounter{section}{3}
\sectiondiv{Dichotomous Keys}{Identifying organisms \textbullet{} Table and flowchart keys \textbullet{} Writing a good question}\label{sec:4}
\lessoncodes{VC2S8U01 \textbullet{} VC2S8I04 \textbullet{} VC2S8I06}

\vocabhead
\begin{ctab}
\begin{tabular}{|SW{3.4cm}|W{7.0cm}|}
  \hline\hd & \thead{Word} & \thead{Definition (copy from the slide)} \\ \hline
  & \scaf{\textbf{Dichotomous key}} & \scaf{A set of two-choice questions that leads to the name of an organism.} \\ \hline
  & \textbf{Identify} & \\ \hline
  & \textbf{Characteristic} & \\ \hline
  & \textbf{Observable} & \\ \hline
  & \textbf{Flowchart} & \\ \hline
\end{tabular}
\end{ctab}

% ─── PART A ───────────────────────────────────────────────────────────────────
\lessonpart{A}{Reading and comprehension}{Teal}\label{read:4}

\begin{readingbox}{Reading --- Twenty Questions for Spiders}
A spider has just run across the shed floor. Is it a harmless orb-weaver or a redback? A field guide with 300 photos is slow to search. A \textbf{dichotomous key} is faster.

``Dichotomous'' means ``cut in two''. A dichotomous key asks one question at a time, and each question has exactly two answers. Each answer sends you either to another question or to a name. After a few steps you \textbf{identify} the organism: you find out what it is.

Every question in a key is about a \textbf{characteristic}, a feature of the organism. A good characteristic is \textbf{observable}: anyone can see or measure it. ``Does it have a red mark on its abdomen?'' is a good question. ``Is it dangerous?'' is not, because you cannot see danger.

Keys come in two forms. A table key lists numbered pairs of statements. A \textbf{flowchart} key draws the questions as boxes joined by arrows labelled with the two answers.
\end{readingbox}

\rcreset
\rcq{What does ``dichotomous'' mean?}{1.5cm}
\rcq{How many answers does each question in a key have, and where can each answer lead?}{2.5cm}
\rcq{What makes a characteristic a good one to use in a key?}{2.5cm}
\rcq{Why is ``Is it dangerous?'' a poor question for a key?}{2.5cm}
\rcq{Name the two forms of dichotomous key.}{2cm}

% ─── PART B ───────────────────────────────────────────────────────────────────
\lessonpart{B}{Theory and examples}{Blue}\label{theory:4}

\subsection{Using a Key}\label{sub:4.1}

\begin{center}
\begin{tikzpicture}[font=\small,
    q/.style={draw=Ink, line width=0.8pt, rounded corners=1.2mm, fill=white, align=center, inner sep=3pt, text width=3.3cm},
    r/.style={draw=Ink, line width=0.6pt, fill=black!8, align=center, inner sep=3pt, text width=2.1cm},
    a/.style={-{Stealth[length=5pt]}, line width=0.8pt, Ink}]
  \node[q] (q1) at (0,0) {Does it have a backbone?};
  \node[q] (q2) at (-2.6,-1.6) {Does it have legs?};
  \node[q] (q3) at (2.6,-1.6) {Does it have six legs?};
  \node[r] (r1) at (-3.9,-3.3) {Eastern grey kangaroo};
  \node[r] (r2) at (-1.3,-3.3) {Eastern brown snake};
  \node[r] (r3) at (1.3,-3.3) {Common grasshopper};
  \node[r] (r4) at (3.9,-3.3) {Garden orb-weaver};
  \draw[a] (q1) -- node[above left, font=\scriptsize]{yes} (q2);
  \draw[a] (q1) -- node[above right, font=\scriptsize]{no} (q3);
  \draw[a] (q2) -- node[left, font=\scriptsize]{yes} (r1);
  \draw[a] (q2) -- node[right, font=\scriptsize]{no} (r2);
  \draw[a] (q3) -- node[left, font=\scriptsize]{yes} (r3);
  \draw[a] (q3) -- node[right, font=\scriptsize]{no} (r4);
\end{tikzpicture}
\figcap{A flowchart key for four Victorian animals.}
\end{center}

\begin{workedex}{Following the key}
Identify an animal with a backbone and no legs.

\smallskip
\textbf{Question 1:} Does it have a backbone? \emph{Yes}, so go left.\\
\textbf{Question 2:} Does it have legs? \emph{No}, so go right.\\
\textbf{Answer:} eastern brown snake.
\end{workedex}

\begin{wedo}{Following the key}
Use the flowchart key above to identify an animal with no backbone and eight legs. Write each question and the answer you chose.
\answerbox{4cm}
\end{wedo}

\subsection{Writing a Good Question}\label{sub:4.2}

\begin{workedex}{Separating a bat from a sugar glider}
\textbf{Poor question:} ``Is it active at night?'' Both are, so it does not separate them.

\smallskip
\textbf{Good question:} ``Are its wings held up by very long finger bones?'' Yes for the bat; no for the sugar glider, whose gliding skin stretches from wrist to ankle. A good question is observable \emph{and} splits the group in two.
\end{workedex}

\begin{wedo}{Writing a good question}
Write one observable question that separates a platypus from an eastern grey kangaroo, and explain why it works.
\answerbox{4cm}
\end{wedo}

\begin{countrybox}{First Nations perspectives --- Identifying with every sense}
Wurundjeri Woi-wurrung knowledge identifies plants and animals with every sense: the smell of a crushed leaf, the call of a frog or bird, texture, taste, and what grows alongside. A dichotomous key uses only what you can see. Both are tools for identifying living things, but they hold different kinds of knowledge.
\end{countrybox}

% ─── PART C ───────────────────────────────────────────────────────────────────
\lessonpart{C}{Practice}{Amber}\label{prac:4}

\task Use this table key to identify each spider.
\begin{ctab}
\begin{tabular}{|C{0.7cm}|L{5.0cm}|L{3.9cm}|}
  \hline\thead{} & \thead{Characteristic} & \thead{Go to / name} \\ \hline
  1a & Body and legs all black & Go to 2 \\ \hline
  1b & Body brown, or with markings & Go to 3 \\ \hline
  2a & Red or orange mark under the abdomen & \textbf{Redback spider} \\ \hline
  2b & No red or orange mark & \textbf{Black house spider} \\ \hline
  3a & Body wider than 3\,cm and hairy & \textbf{Huntsman spider} \\ \hline
  3b & Body narrower than 3\,cm and smooth & \textbf{Garden orb-weaver} \\ \hline
\end{tabular}
\end{ctab}
\pt{a} A small black spider with a red stripe under its abdomen.
\answerbox{1cm}
\pt{b} A large, hairy brown spider wider than your palm.
\answerbox{1cm}
\pt{c} A black spider with no coloured marks anywhere.
\answerbox{1cm}
\pt{d} A small, smooth brown spider with a pale pattern.
\answerbox{1cm}

\needspace{14\baselineskip}
\task Design a table key to identify four Victorian birds: little penguin (cannot fly, black and white), superb lyrebird (very long showy tail, brown), wedge-tailed eagle (huge, dark, hooked beak, wedge-shaped tail) and laughing kookaburra (large head, thick straight beak, brown and white). Use features you can see, not behaviour.
\begin{ctab}
\begin{tabular}{|SZ{0.9cm}|W{5.0cm}|W{3.9cm}|}
  \hline\hd & \thead{} & \thead{Characteristic} & \thead{Go to / name} \\ \hline
  & \scaf{1a} & \scaf{Wings used for swimming, cannot fly} & \scaf{Little penguin} \\ \hline
  & 1b & & \\ \hline
  & 2a & & \\ \hline
  & 2b & & \\ \hline
  & 3a & & \\ \hline
  & 3b & & \\ \hline
\end{tabular}
\end{ctab}

\needspace{10\baselineskip}
\task Give one advantage and one disadvantage of a flowchart key compared with a table key.
\answerbox{5cm}

% ─── PART D ───────────────────────────────────────────────────────────────────
\lessonpart{D}{Above level}{Violet}\label{ext:4}

\begin{yr8box}
These questions go beyond the Year 7 standard.

\medskip
\begin{keepext}
\exttask Keys go wrong when members of one species look different: young and adults, or males and females. Describe how this could make a key give the wrong answer, and suggest how to improve the key.
\answerbox{5cm}
\end{keepext}
\begin{keepext}
\exttask DNA barcoding identifies a species from a short piece of its DNA. A scientist finds an unknown beetle in a Victorian forest. Give one advantage and one limitation of DNA barcoding compared with a dichotomous key.
\answerbox{5cm}
\end{keepext}
\end{yr8box}
